\documentclass[11pt]{article}
\usepackage[utf8]{inputenc}
\usepackage[T1]{fontenc}
\usepackage{amsmath,amssymb,amsthm}
\usepackage{geometry}
\usepackage{hyperref}
\usepackage{tcolorbox}
\usepackage{booktabs}
\usepackage{enumitem}
\geometry{margin=1in}
\tcbuselibrary{breakable, skins}

% Hyperlink colors
\hypersetup{
  colorlinks=true,
  linkcolor=blue!50!black,
  urlcolor=blue!50!black,
  citecolor=blue!50!black
}

% Theorem environments
\newtheorem{theorem}{Theorem}[section]
\newtheorem{proposition}[theorem]{Proposition}
\newtheorem{lemma}[theorem]{Lemma}
\newtheorem{corollary}[theorem]{Corollary}
\theoremstyle{definition}
\newtheorem{definition}[theorem]{Definition}
\theoremstyle{remark}
\newtheorem*{remark}{Remark}

% Custom commands
\newcommand{\E}{\mathbb{E}}
\newcommand{\Var}{\mathrm{Var}}
\newcommand{\Cov}{\mathrm{Cov}}
\renewcommand{\Pr}{\mathbb{P}}
\newcommand{\N}{\mathbb{N}}
\newcommand{\R}{\mathbb{R}}
\newcommand{\eps}{\varepsilon}
\newcommand{\half}{\tfrac{1}{2}}

% Intuition box
\newtcolorbox{intuition}{
  colback=blue!5!white,
  colframe=blue!50!black,
  fonttitle=\bfseries,
  title=Intuition,
  breakable
}

\title{\textbf{Porter-Thomas Statistics under Two Sampling Regimes}\\[0.5em]
\large Why ``many shots, one circuit'' and ``many circuits, one shot'' \\ 
give the same expected F\textsubscript{XEB} but different concentration\\[0.5em]
\large \textit{A side-by-side variance analysis}}
\author{Quip Protocol Project --- Internal Reference Document}
\date{June 2026}

\begin{document}

\maketitle

\begin{abstract}
\noindent
Both the Aaronson--Hung (AH) and Liu et al.\ protocols sample from random
quantum circuits and use the cross-entropy benchmark F\textsubscript{XEB} as
their acceptance test. They differ in \textit{how} the samples are collected:
AH uses one circuit with $m$ shots per round; Liu uses $m$ distinct circuits
with one shot each. This document establishes that both regimes yield
identically Porter-Thomas-distributed observed probabilities, produce the
same expected F\textsubscript{XEB}, but differ in their variance structure
in a way that has direct consequences for certified-entropy bounds. The
core result is a variance decomposition (\S\ref{sec:variance})
demonstrating that AH's samples are correlated through the shared circuit
while Liu's samples are independent across rounds. This correlation gap
explains why AH requires the Entropy Accumulation Theorem for its
concentration argument and gives a looser bound, while Liu's i.i.d.\ structure
admits direct Chernoff concentration and a tighter certified-entropy claim
for the same hardware. Quip's v0.4 implementation adopts Liu's structure.
\end{abstract}

\tableofcontents
\vspace{1em}

\section{Setup and Notation}

We consider an $n$-qubit quantum circuit $C$ drawn from a distribution
sufficiently close to the Haar measure on $U(2^n)$. Throughout this
document, write $N := 2^n$ for the dimension of the computational basis.
For any circuit $C$ and bitstring $x \in \{0,1\}^n$, define
\[
  p_C(x) := |\langle x \mid U_C \mid 0^n \rangle|^2,
\]
the probability that an ideal noiseless execution of $C$ produces measurement
outcome $x$.

A real (noisy) quantum sampler with fidelity $\phi \in [0,1]$ is modelled
as producing $x$ from the mixture distribution
\[
  q_C(x) = \phi \cdot p_C(x) + (1-\phi) \cdot \frac{1}{N},
\]
i.e., the ideal distribution with probability $\phi$ and uniform noise with
probability $1-\phi$.

\subsection{The two sampling regimes}

We distinguish:

\begin{description}
\item[Regime A (Aaronson--Hung):] Pick one circuit $C$. Take $m$ samples
  $x_1, x_2, \ldots, x_m$ i.i.d.\ from $q_C$. The $m$ samples share the same
  underlying circuit.
\item[Regime B (Liu et al.):] Pick $m$ independent circuits $C_1, C_2,
  \ldots, C_m$. From each $C_i$, take exactly one sample $x_i \sim
  q_{C_i}$. Each $(C_i, x_i)$ pair is independent of every other.
\end{description}

The cross-entropy benchmarking score is identical in form:
\begin{equation}
  F_{\mathrm{XEB}} := \frac{N}{m} \sum_{i=1}^{m} p_{C_i}(x_i) - 1,
  \label{eq:fxeb}
\end{equation}
where in Regime A, all $C_i = C$ (the shared circuit), and in Regime B
each $C_i$ is independent. The protocol accepts if $F_{\mathrm{XEB}} \geq
\chi$ for some threshold $\chi$ (Liu uses $\chi = 0.3$).

\subsection{What we want to establish}

Two questions:
\begin{enumerate}[label=(\arabic*)]
\item Is the \emph{statistical distribution} of $p_{C_i}(x_i)$ identical in
  the two regimes? \textbf{Yes.}
\item Is the \emph{concentration behaviour} of $F_{\mathrm{XEB}}$ identical
  in the two regimes? \textbf{No.}
\end{enumerate}

The first answer is the source of common confusion --- both regimes
produce Porter-Thomas statistics, so it's tempting to conclude they're
interchangeable. The second answer is why they're not, and why Liu's
structure is preferred in v0.4.

\section{The Porter-Thomas distribution}
\label{sec:pt}

We collect the relevant facts about $p_C(x)$ for Haar-random circuits.

\begin{theorem}[Porter--Thomas]
\label{thm:pt}
Let $C$ be a circuit whose unitary is Haar-distributed on $U(N)$. For any
fixed bitstring $x \in \{0,1\}^n$, the random variable
$N \cdot p_C(x)$ converges in distribution to $\mathrm{Exp}(1)$ as
$N \to \infty$.
\end{theorem}

The proof appears in Porter and Thomas (1956) and is reviewed in
Boixo et al.\ (2018). Sufficiently-deep random circuits approximate the
Haar distribution closely enough that the convergence holds in practice
once $d \gtrsim n$ (where $d$ is the circuit depth).

\subsection{Moments of Porter-Thomas}

For $Y \sim \mathrm{Exp}(1)$, the moment generating function is
$\E[e^{tY}] = 1/(1-t)$ for $t < 1$, giving moments $\E[Y^k] = k!$. Since
$N \cdot p_C(x) \sim Y$, we have
\begin{equation}
  \E[p_C(x)^k] = \frac{k!}{N^k}.
  \label{eq:moments}
\end{equation}
The first three moments will appear throughout:
\[
  \E[p_C(x)] = \frac{1}{N}, \qquad
  \E[p_C(x)^2] = \frac{2}{N^2}, \qquad
  \E[p_C(x)^3] = \frac{6}{N^3}.
\]

\subsection{The collision probability identity}

A computation we'll use:
\begin{lemma}
\label{lem:collision}
For Haar-random $C$,
\[
  \E_C\!\left[\sum_{x \in \{0,1\}^n} p_C(x)^2 \right] = \frac{2}{N}.
\]
\end{lemma}

\begin{proof}
By linearity and \eqref{eq:moments}, $\E_C[\sum_x p_C(x)^2] = N \cdot
\E_C[p_C(x)^2] = N \cdot 2/N^2 = 2/N$. \qedhere
\end{proof}

This identity is what produces the famous ``factor of 2'' in F\textsubscript{XEB}
calculations.

\section{The size-biased distribution: $p_C(x)$ when $x \sim p_C$}
\label{sec:sizebias}

Both regimes involve a quantity of the form $p_{C_i}(x_i)$ where $x_i$
\emph{was actually drawn from} $q_{C_i}$ (not a fixed bitstring chosen
independently of the circuit). This conditional choice changes the
distribution.

\begin{intuition}
If you sample from a distribution, you're more likely to observe bitstrings
the distribution assigned high probability to. So the values of $p_C(x)$
you see when $x \sim p_C$ are biased toward larger values than $p_C(x)$
for a fixed independent $x$. This is called \emph{size-biased sampling}.
\end{intuition}

\subsection{Conditional distribution: ideal sampler}

For the noiseless case $\phi = 1$ (so $q_C = p_C$):

\begin{proposition}
\label{prop:sizebias}
Let $C$ be Haar-random and $X \sim p_C$. Then for any measurable bounded
$f$,
\[
  \E[f(p_C(X))] = N \cdot \E[p_C(x_0)\, f(p_C(x_0))]
\]
where $x_0$ is any fixed bitstring. In particular:
\begin{align*}
  \E[p_C(X)] &= \frac{2}{N}, \\
  \E[p_C(X)^2] &= \frac{6}{N^2}, \\
  \mathrm{Var}[p_C(X)] &= \frac{6}{N^2} - \frac{4}{N^2} = \frac{2}{N^2}.
\end{align*}
\end{proposition}

\begin{proof}
For any function $f$ of the probability,
\begin{align*}
\E_{C, X \sim p_C}[f(p_C(X))]
  &= \E_C\!\left[\sum_x p_C(x) f(p_C(x))\right] \\
  &= N \cdot \E_C[p_C(x_0) f(p_C(x_0))]
\end{align*}
where the second line uses symmetry over $x$ (Haar-randomness makes all
bitstrings statistically equivalent). Setting $f(y) = y^k$ and applying
\eqref{eq:moments} gives
\[
  \E[p_C(X)^k] = N \cdot \E[p_C(x_0)^{k+1}] = N \cdot \frac{(k+1)!}{N^{k+1}}
  = \frac{(k+1)!}{N^k}.
\]
Setting $k=1, 2$ gives the stated values. \qedhere
\end{proof}

\begin{intuition}
The conditional distribution of $N \cdot p_C(X)$ when $X \sim p_C$
is the \emph{Gamma}$(2,1)$ distribution: a sum of two independent
$\mathrm{Exp}(1)$'s, with mean 2. Compare to $N \cdot p_C(x_0)$ for fixed
$x_0$, which is $\mathrm{Exp}(1)$ with mean 1. The size-biased version is
literally ``one extra exponential's worth'' of probability mass.
\end{intuition}

\subsection{Conditional distribution: noisy sampler}

For $\phi < 1$, $X \sim q_C = \phi p_C + (1-\phi)/N$, so by linearity:
\begin{equation}
  \E[p_C(X) \mid C] = \phi \cdot \E_{X \sim p_C}[p_C(X) \mid C] + (1-\phi)
    \cdot \E_{X \sim U}[p_C(X) \mid C].
  \label{eq:noisycond}
\end{equation}
The first term gives $\phi$ times the size-biased mean; the second term
gives $(1-\phi)/N$. Averaging over $C$:
\begin{equation}
  \E[p_C(X)] = \phi \cdot \frac{2}{N} + (1-\phi) \cdot \frac{1}{N}
   = \frac{1 + \phi}{N}.
  \label{eq:noisyexp}
\end{equation}

This is the key identity: \textbf{the expected probability of an observed
sample under a fidelity-$\phi$ sampler is $(1+\phi)/N$.}

\section{Expected F\textsubscript{XEB} --- identical in both regimes}
\label{sec:exp}

We can now compute $\E[F_{\mathrm{XEB}}]$ in both regimes.

\subsection{Regime B (Liu)}

In Regime B, each $(C_i, x_i)$ is an independent draw. Using
\eqref{eq:noisyexp}:
\begin{equation}
  \E[F_{\mathrm{XEB}}^{(B)}]
   = \frac{N}{m} \sum_{i=1}^m \E[p_{C_i}(x_i)] - 1
   = N \cdot \frac{1+\phi}{N} - 1
   = \phi.
   \label{eq:expB}
\end{equation}

\subsection{Regime A (AH)}

In Regime A, all $m$ samples come from the same $C$. We need the law of
total expectation:
\begin{align}
  \E[F_{\mathrm{XEB}}^{(A)}]
   &= \E_C\!\left[ \frac{N}{m} \sum_{i=1}^m \E[p_C(x_i) \mid C] - 1
   \right] \notag \\
   &= \E_C\!\left[ \frac{N}{m} \cdot m \cdot \E[p_C(x_1) \mid C] -1 \right]
   \notag \\
   &= N \cdot \E_C\!\left[ \E[p_C(x_1) \mid C] \right] - 1 \notag \\
   &= N \cdot \E[p_C(x_1)] - 1 \notag \\
   &= N \cdot \frac{1+\phi}{N} - 1 = \phi.
   \label{eq:expA}
\end{align}

\begin{intuition}
Both regimes give $\E[F_{\mathrm{XEB}}] = \phi$. This is what makes
F\textsubscript{XEB} a useful fidelity estimator in both settings. The
\emph{mean} is structure-agnostic --- it depends only on the per-sample
expectation, which involves Porter-Thomas size-biasing in the same way for
both regimes.
\end{intuition}

\begin{remark}
For an honest quantum sampler with $\phi \approx 1$, $\E[F_{\mathrm{XEB}}]
\approx 1$. For a uniform-random classical adversary with $\phi = 0$,
$\E[F_{\mathrm{XEB}}] = 0$. The Liu protocol's threshold $\chi = 0.3$ sits
between these, corresponding to roughly the fidelity Quantinuum H2 delivers
on the 56-qubit circuit ensemble.
\end{remark}

\section{Variance --- where the regimes differ}
\label{sec:variance}

Now we come to the central calculation. We compute $\Var[F_{\mathrm{XEB}}]$
in each regime.

\subsection{Variance in Regime B (Liu)}

Each term $p_{C_i}(x_i)$ in the sum is an independent random variable
(independent across $i$ because the $(C_i, x_i)$ pairs are independent).
So
\begin{equation}
  \Var[F_{\mathrm{XEB}}^{(B)}] = \frac{N^2}{m^2} \cdot \sum_{i=1}^m
  \Var[p_{C_i}(x_i)] = \frac{N^2}{m} \cdot \Var[p_C(X)],
  \label{eq:varB}
\end{equation}
where in the last step we use that the $m$ variances are equal.

We computed $\Var[p_C(X)] = 2/N^2$ for the noiseless case (Proposition
\ref{prop:sizebias}). For $\phi < 1$ the variance changes by an $O(1)$
factor; what matters is that
\begin{equation}
  \Var[F_{\mathrm{XEB}}^{(B)}] = \frac{c(\phi)}{m}
  \label{eq:varB-rate}
\end{equation}
for some constant $c(\phi) = O(1)$. \textbf{Variance scales as $1/m$.}

\subsection{Variance in Regime A (AH)}

For Regime A we use the \emph{law of total variance}:
\begin{equation}
  \Var[F_{\mathrm{XEB}}^{(A)}] = \E_C\big[\Var[F_{\mathrm{XEB}}^{(A)} \mid C]\big]
  + \Var_C\big[\E[F_{\mathrm{XEB}}^{(A)} \mid C]\big].
  \label{eq:lotv}
\end{equation}

We compute each term.

\subsubsection*{First term: within-circuit variance}

Given $C$ fixed, the samples $x_1, \ldots, x_m$ are i.i.d.\ from $q_C$. So
\begin{equation}
  \Var[F_{\mathrm{XEB}}^{(A)} \mid C] = \frac{N^2}{m^2} \cdot m \cdot
  \Var[p_C(X) \mid C] = \frac{N^2}{m} \cdot \Var[p_C(X) \mid C].
\end{equation}
Taking the outer expectation:
\begin{equation}
  \E_C\big[\Var[F_{\mathrm{XEB}}^{(A)} \mid C]\big]
  = \frac{N^2}{m} \cdot \E_C[\Var[p_C(X) \mid C]].
  \label{eq:within}
\end{equation}

This is the term that scales as $1/m$. So far, indistinguishable from
Regime B.

\subsubsection*{Second term: between-circuit variance --- the new term}

The conditional expectation given $C$ is $\E[F_{\mathrm{XEB}}^{(A)} \mid C]
= N \cdot \E[p_C(X) \mid C] - 1$. This is itself a random variable
(depending on which $C$ was drawn).

For the noiseless case, $\E[p_C(X) \mid C] = \sum_x p_C(x)^2$ (this is the
\emph{collision probability} of distribution $p_C$). By
Lemma~\ref{lem:collision}, its expectation is $2/N$. Its variance is
non-trivial:
\begin{lemma}
\label{lem:between}
For Haar-random $C$,
\[
  \Var_C\!\left[\sum_x p_C(x)^2\right] = \frac{20}{N^3} + O(N^{-4}).
\]
\end{lemma}

(The leading-order computation involves $\E_C[(\sum_x p_C(x)^2)^2] = 4/N^2
+ 20/N^3 + O(N^{-4})$, which follows from Weingarten calculus on the
Haar measure. We don't reproduce the calculation here; see the Boixo et
al.\ supplement.)

So:
\begin{align}
  \Var_C\big[\E[F_{\mathrm{XEB}}^{(A)} \mid C]\big]
   &= \Var_C\big[N \cdot \E[p_C(X) \mid C] - 1\big] \notag\\
   &= N^2 \cdot \Var_C[\E[p_C(X) \mid C]] \notag \\
   &\approx N^2 \cdot \frac{20}{N^3} \notag \\
   &= \frac{20}{N}.
   \label{eq:between}
\end{align}

\textbf{Crucially, this term does not depend on $m$.}

\subsubsection*{Putting it together}

Combining \eqref{eq:within} and \eqref{eq:between}:
\begin{equation}
  \Var[F_{\mathrm{XEB}}^{(A)}] = \underbrace{\frac{c(\phi)}{m}}_{\text{within-circuit}}
   + \underbrace{\frac{20}{N} + O(\cdot)}_{\text{between-circuit}}.
   \label{eq:varA}
\end{equation}

\subsection{Comparison}

Compare:
\begin{align*}
  \Var[F_{\mathrm{XEB}}^{(B)}] &= \frac{c(\phi)}{m}, \\
  \Var[F_{\mathrm{XEB}}^{(A)}] &= \frac{c(\phi)}{m} + \frac{20}{N} +
  O(\cdot).
\end{align*}

For $n = 56$, $N = 2^{56} \approx 7.2 \times 10^{16}$, so $20/N \approx
3 \times 10^{-16}$. For $m = 1522$, $c(\phi)/m \sim 1/m \approx 6.6
\times 10^{-4}$. So the between-circuit term is dwarfed by the
within-circuit term.

\textit{Does this mean the difference is negligible?} No --- and this is
the subtle point. The variance values \emph{evaluated for a single
random circuit} are similar in magnitude. But the structural difference
manifests when we ask about \emph{concentration over many independent
protocol executions}.

\begin{intuition}
The two regimes give similar single-run variance because the
between-circuit fluctuation is small ($N$ is huge, so individual circuits
all look approximately Porter-Thomas). The structural difference is not
in the magnitude of variance but in \textbf{which concentration
inequality applies}: i.i.d.\ samples (Regime B) admit standard Chernoff
bounds; correlated samples sharing a circuit (Regime A) require the
Entropy Accumulation Theorem.
\end{intuition}

\section{Concentration: where the structural difference bites}
\label{sec:concentration}

The certified-entropy bound depends not just on the variance of
F\textsubscript{XEB} but on \emph{tail bounds} --- the probability that
F\textsubscript{XEB} deviates from its expectation by more than $\delta$,
for various $\delta$.

\subsection{Regime B: direct Chernoff}

In Regime B, the $m$ summands $p_{C_i}(x_i)$ are i.i.d.\ bounded random
variables (bounded by $1$, since they're probabilities). The Chernoff
bound (or Hoeffding's inequality) gives
\begin{equation}
  \Pr[|F_{\mathrm{XEB}}^{(B)} - \phi| > \delta] \leq 2 \exp(-c \cdot m
  \delta^2),
\end{equation}
for some constant $c$. The protocol's soundness analysis uses
exactly this bound (with a careful split of the soundness budget
$\eps_\mathrm{sou}/2$ between this Chernoff term and the hypergeometric
tail; see Liu et al.\ Lemma 2 and our own implementation in
\texttt{leg5\_qmin\_liu.compute\_q\_min}).

\textbf{Crucially: nothing more sophisticated is needed.} The bound is
sharp, computable in closed form, and yields the tight $Q_\mathrm{min}$
calibration.

\subsection{Regime A: Entropy Accumulation Theorem}

In Regime A, the $m$ summands $p_C(x_i)$ share the single random variable
$C$. They are not independent; they're conditionally i.i.d.\ given $C$.

A direct Chernoff bound \emph{does not work}: the variance includes a
between-circuit term that's not reduced by increasing $m$.

Aaronson and Hung handle this with the Entropy Accumulation Theorem (EAT)
of Dupuis, Fawzi, and Renner (2020). EAT is a chain-rule-style decomposition
that handles the conditional structure correctly. The high-level statement:
\begin{quote}
Even when later rounds are conditioned on earlier rounds' outcomes, the
total entropy accumulates linearly (modulo correction terms scaling with
$\sqrt{m}$) provided per-round entropy is appropriately characterised.
\end{quote}

For RCS with $m$ shots on a single circuit, EAT gives a certified-entropy
bound of the form
\begin{equation}
  H_{\min}^{\mathrm{cert}} \geq m \cdot h(\phi) - O(\sqrt{m}) -
  \log(1/\eps_\mathrm{sou})
\end{equation}
for some per-shot entropy rate $h(\phi)$.

\subsection{Why Liu's bound is tighter}

The crucial comparison: AH's bound has an $O(\sqrt{m})$ correction term;
Liu's i.i.d.\ Chernoff has an $O(\log(1/\eps)/\sqrt{m})$ correction in the
deviation.

For typical parameters ($m \sim 10^3$, $\eps_\mathrm{sou} \sim 10^{-6}$),
the $O(\sqrt{m}) \approx 30$ correction in AH \textbf{eats into the
certified-entropy budget} more than the $O(\log(1/\eps))$ correction in
Liu's i.i.d.\ regime.

This is why, for the same hardware and the same nominal $F_{\mathrm{XEB}}$,
Liu's structure yields \emph{more} certified bits per round than AH's
structure does. \textbf{Same physics; same expectation; tighter bound from
the i.i.d.\ structure.}

\section{Numerical comparison for Liu's parameter regime}

For $n=56$, $m=1522$, $M = 30{,}010$, $\chi = 0.3$, $\eps_\mathrm{sou} =
10^{-6}$, and the calibrated Liu adversary model ($4 \times$ Frontier
theoretical, $\mathrm{eff} = 0.5$):

\begin{center}
\begin{tabular}{lrr}
\toprule
\textbf{Quantity} & \textbf{Regime A (AH)} & \textbf{Regime B (Liu)} \\
\midrule
$\E[F_{\mathrm{XEB}}]$ & $\phi$ (same) & $\phi$ (same) \\
$\Var[F_{\mathrm{XEB}}]$ & $\sim c(\phi)/m$ & $\sim c(\phi)/m$ \\
Concentration tool & EAT & i.i.d.\ Chernoff \\
Soundness-eating correction & $O(\sqrt m)$ & $O(\log(1/\eps))$ \\
Certified bits per round & smaller & larger \\
$Q_\mathrm{min}$ (Liu's parameters) & $\approx $ larger$^\dagger$ & $1297$ \\
$H_{\min}^{\mathrm{cert}}$ & smaller$^\dagger$ & $71{,}313$ bits \\
\bottomrule
\end{tabular}
\end{center}

\noindent {\small $^\dagger$ ``Larger $Q_\mathrm{min}$'' for AH means the
threshold is more conservative --- it certifies less for the same
acceptance. Exact AH numbers depend on the EAT-derived bound; the
qualitative point is that AH is strictly looser for the same parameters.}

\section{Implications for Quip}

Three concrete implications for the protocol design:

\paragraph{1. Liu's structure is the right choice.} For the same hardware
and the same $F_{\mathrm{XEB}}$ score, Liu's i.i.d.\ structure certifies more
bits per round than AH's. The v0.4 refactor adopted this structure, and
this document explains \emph{why} that adoption was the right call.

\paragraph{2. F\textsubscript{XEB} computation is structure-agnostic.}
The formula \eqref{eq:fxeb} is identical in both regimes. The function in
\texttt{leg5\_certified\_entropy.compute\_fxeb} operates on (sample,
probability) pairs without caring how the samples were grouped. This is
why verifying our F\textsubscript{XEB} implementation against Liu's
experimental data is a meaningful test of the formula, independent of the
structural choice.

\paragraph{3. The structural choice is separable from hardware.} Whether
we run on Quantinuum H2, IonQ Forte, or another QPU, the
expected $F_{\mathrm{XEB}}$ is set by hardware fidelity. The structural
choice (AH vs Liu) is independent: it's a protocol-level decision about
how to allocate circuits to rounds. Hardware decisions and structural
decisions can be made independently.

\section{Summary}

The two questions and their answers:

\begin{enumerate}
\item \textbf{Is the statistical distribution of $p_{C_i}(x_i)$
identical in the two regimes?} \\
\textbf{Yes.} In both regimes, $p_{C_i}(x_i)$ follows the size-biased
Porter-Thomas distribution, with mean $(1+\phi)/N$.
\item \textbf{Is the concentration behaviour of F\textsubscript{XEB}
identical in the two regimes?} \\
\textbf{No.} Regime B's samples are independent, admitting direct Chernoff
concentration. Regime A's samples are correlated through the shared
circuit, requiring the Entropy Accumulation Theorem and producing a
strictly looser certified-entropy bound.
\end{enumerate}

The expected $F_{\mathrm{XEB}}$ is identical; the variance has the same
leading $1/m$ scaling; but the \emph{tail behaviour} differs in a way
that propagates through the certified-entropy chain. Liu's structure is
optimal among the two; this is what v0.4 implements.

\section*{References}

\begin{description}
\item[Aaronson, S.\ and Hung, S.-H.] (2023).
  \emph{Certified Randomness from Quantum Supremacy}. STOC.

\item[Boixo, S.\ et al.] (2018).
  \emph{Characterizing quantum supremacy in near-term devices}. Nature
  Physics 14, 595--600.

\item[Dupuis, F., Fawzi, O., Renner, R.] (2020).
  \emph{Entropy Accumulation}. Communications in Mathematical Physics 379,
  867--913.

\item[Liu, M.\ et al.] (2025).
  \emph{Certified randomness using a trapped-ion quantum processor}.
  Nature 640:343; arXiv:2503.20498.

\item[Porter, C.E.\ and Thomas, R.G.] (1956).
  \emph{Fluctuations of nuclear reaction widths}. Physical Review 104,
  483.
\end{description}

\vfill

\noindent\rule{\linewidth}{0.4pt}\\[2pt]
\textit{Quip Protocol Project --- Internal reference. Last revised June 2026.
Companion documents: \texttt{porter\_thomas\_report.pdf} (full
Porter-Thomas derivations), \texttt{leg5\_verification\_design.pdf}
(verification deck), and \texttt{quip\_design\_notes.pdf} (consolidated
design notes).}

\end{document}
